% Phase 9 independent adversarial result; self-contained mathematical insert.
\section{A missing hypothesis in an unequal-weight three-point reduction}

For weighted scalar locations, write $P_{UV}$ for the unnormalized
two-point Jensen distortion and $T_{UVW}$ for the three-point distortion.
The three original weights are retained in each expression.

\begin{proposition}[The compatible-optimizer reduction survives weights]
For four distinct ordered scalar locations with arbitrary positive source
weights and a differentiable strictly convex generator, there are
compatible optimal two- and three-cell partitions unless an optimal
three-cell partition merges the middle pair $BC$ and the only optimal
contiguous two-cell partition is $AB\mid CD$.
\end{proposition}
\begin{proof}
Flat optima can be chosen contiguous: for fixed ordered centers, the
difference of two scalar Bregman costs is affine in the data location, so
assignment to a minimizing center has ordered interval cells; recomputing
the centers can only improve the objective. Positive point weights do not
change which center minimizes the cost at a given point.

For every ordered triple, $T_{ABC}\geq P_{AB}+P_{BC}$. One way to see
this is through curvature kernels. Adding $C$ to the cell $AB$ adds the
nonnegative Jensen kernel for merging its weighted mean with $C$;
therefore $K_{ABC}\geq K_{AB}$. Likewise $K_{ABC}\geq K_{BC}$, and the
two latter kernels have disjoint open supports. Equivalently the same
argument works by the distributional curvature measure of a convex
generator. Thus their sum is pointwise dominated by $K_{ABC}$.

If an optimal adjacent fine pair is $AB$, the only incompatible
contiguous coarse split is $A\mid BCD$. Its cost obeys
$T_{BCD}\geq P_{BC}+P_{CD}\geq P_{AB}+P_{CD}$, so the compatible split
$AB\mid CD$ is at least as good. The $CD$ case is symmetric.
If the optimal pair is $BC$, the coarse splits $A\mid BCD$ and
$ABC\mid D$ are already compatible; only $AB\mid CD$ can conflict.
\end{proof}

\begin{proposition}[The two three-point inequalities alone give no bound]
Let $q=P_{BC}$, $S=P_{AB}+P_{BC}$, and $T=T_{ABC}$. There is no finite
degree-only bound on a parameter $R$ from the inequalities
\[
 q\leq T/R^2,\qquad S\leq2T/R
\]
for arbitrary positive triple weights, even if the curvature is the
constant polynomial $g\equiv1$ and the locations are fixed.
\end{proposition}
\begin{proof}
Take $\varphi(x)=x^2/2$, $(A,B,C)=(0,1/2,1)$, and original weights
$(1,\varepsilon^2,\varepsilon)$ for $0<\varepsilon\leq1$.
Common normalization of all weights is immaterial. The exact weighted
variance identities give
\[
 P_{AB}=\frac{\varepsilon^2}{8(1+\varepsilon^2)},\qquad
 q=\frac{\varepsilon^2}{8(1+\varepsilon)},\qquad
 T=\frac{\varepsilon+(\varepsilon^2+\varepsilon^3)/4}
          {2(1+\varepsilon+\varepsilon^2)}.
\]
In particular $q\leq\varepsilon^2/8$, $S\leq\varepsilon^2/4$, and
$T\geq\varepsilon/6$. Put $R=\varepsilon^{-1/2}$. Then
\[
 q\leq\varepsilon^2/8\leq\varepsilon T=T/R^2,
 \qquad
 S\leq\varepsilon^2/4\leq\varepsilon^{3/2}/3
    \leq2\sqrt\varepsilon\,T=2T/R.
\]
As $\varepsilon\downarrow0$, $R\to\infty$ while the degree remains
zero and all locations remain fixed.
\end{proof}

\paragraph{Consequence for a proposed sharp weighted upper bound.}
The compatible-optimizer reduction does not require equal weights.
The localization step of the equal-mass proof does require them, or a
replacement hypothesis: the displayed two inequalities for one triple
are insufficient. A valid general weighted proof must also exploit
information from the fourth point, for example simultaneous control of
both triples relative to the same outer-pair sum, and must check every
dependence on the ratios of the three source masses. This counterexample
does not challenge the Phase 8 quadratic upper bound or its established
weighted lower family.
